
(\(c_i\) \uline{et} \(c_i^{(j)}\)
\uline{étant effectués du poids} \(i\)), \uline{et ces} \(n\) \uline{relations engendrent l'idéal des relations entre les} \(c_i^{(j)}\).
\end{enumerate}

\textbf{Corollaire 4.4.} \uline{Les éléments de} \(C\) \uline{de la forme}~:
\[
  \xi_1^{\alpha_1}\ldots\xi_{n-1}^{\alpha_{n-1}},
  \qquad 0\le\alpha_i\le n-i,
\]
\uline{forment une base de} \(C\) \uline{considéré comme module sur} \(A\).

Pour la démonstration du lemme et du corollaire, nous renvoyons à [2], exposé 4, pages 4.19, auquel ils sont empruntés.

4.5.
Soient \(S\) un schéma quasi-compact, \(\underline{E}\) un \(\underline{O}_S\)-Module localement libre de rang \(n\), \(\pi=(p_1,\ldots,p_k)\) une suite d'entiers \(\geq1\), de somme \(n\). Posons \(X=\underline{\underline{D}}_\pi(\underline{E})\) ; \(X\) est muni d'une suite de quotients canonique \(\underline{L}_i\) ; on pose \(\underline{L}_0=0\), \(\underline{L}_k=f^*(\underline{E})\), \(f\) étant le morphisme structural \(X\to S\). Pour \(1\le i\le k\) on note \(\underline{F}_i\) le noyau de \(\underline{L}_i\to \underline{L}_{i-1}\). Soit \(K^{\bullet}(X)\) le groupe de Grothendieck des faisceaux localement libres de type fini sur \(X\) (resp. des complexes parfaits). On note \(F_j\) la classe de \(\underline{F}_j\), et \(\lambda_i^{(j)}\) la classe de \(\wedge^i(\underline{F}_j)\). Dans le cas où \(K^{\bullet}(X)\) est défini en termes de complexes parfaits, on donne encore un sens évident aux séries \(\lambda_t(E)\), \(\lambda_t(F_j)\), où \(E\) est la classe de \(\underline{E}\). Enfin, on remarque que~:
\[
  \sum_{j=1}^{k}F_j=E.
\]

\textbf{Théorème 4.6.} \uline{Avec les hypothèses et les notations de} 4.5, \uline{la} \(K^{\bullet}(S)\)-\uline{algèbre} \(K^{\bullet}(X)\) \uline{est engendrée par les} \(\lambda_i^{(j)}\); \uline{l'idéal des relations entre les} \(\lambda_i^{(j)}\) \uline{est engendré par les relations déduites de}~:
\[
  \lambda_t(E)=\prod_{j=1}^{k}\lambda_t(F_j).
\]
